% !TEX root = main.tex
\documentclass[main.tex]{subfiles}
\begin{document}

% DRAFT (rewrite of Section 5 following sec5_rewrite_plan.md).
% Previous versions: flow_v1.tex (implicit / semi-implicit), flow_v2.tex, flow_v3.tex (before this rewrite).

\section{Power Particles Revisited}
\label{sec:flow}

In this section, we apply the constrained power diagrams of \S\ref{sec:contributions} to the discretization of incompressible flows, in which each fluid parcel is a cell of a power diagram with prescribed volume, represented by a particle at the centroid of its cell.
We first show that this discretization leads to a holonomic constraint, for which we propose a variational but costly solver (\S\ref{sec:flow-holonomic}).
We then propose two relaxations: a penalty, which keeps the integrator variational but approximates the constraint (\S\ref{sec:flow-softened}), and a projection, which keeps the constraint exact but approximates the integrator (\S\ref{sec:flow-projection}).
Finally, we show that the variational scheme of \citet{Gallouet2018Lagrangian} and Power Particles \cite{deGoes2015PowerParticles} are special cases and simplifications of these relaxations (\S\ref{sec:flow-prior}).

\subsection{Holonomic Formulation}
\label{sec:flow-holonomic}

\paragraph{Spatial discretization.}
For simplicity, we assume a uniform density of one, and we discretize the fluid into $n$ parcels, each one represented by a cell $\cell_i$ of prescribed volume $\pvol_i$, and thus of mass $\pvol_i$, and by a particle $\qq_i$ located at the centroid of its cell.
Building upon \S\ref{sec:joint-constraints}, we set both the sites $\qq$ and the prescribed centroids $\tcentroid$ of $\jointdual$ to the particles, \ie $\tcentroid = \qq$, and define the potential
\begin{equation}
  \cpot_\infty(\qq) = \max_{\shift, \weight} \jointdual(\shift, \weight).
  \label{eq:flow-centroid-potential}
\end{equation}
When the maximum is attained, its maximizer generates a power diagram with volumes $\pvol$ and centroids $\qq$, and $\cpot_\infty$ equals the quantization cost $\smash{\sum_i \int_{\cell_i} \norm{\xx - \qq_i}^2 d\xx}$ of this partition.
However, as discussed in \S\ref{sec:remarks}, the maximum is attained only when the particles are realizable.
The particles are therefore subject to a \emph{holonomic} constraint, \ie a constraint on the positions alone, independent of the velocities, which requires them to remain realizable.
Since realizable centroids lie on the relative boundary of the achievable centroids $\achievable$, with the generators $\dual$ as outward normal (\S\ref{sec:background}), this is a single scalar constraint within the affine subspace \eqref{eq:partition-feasibility-centroid}, whose normal is $\dual$.

Without loss of generality, we place the origin at the centroid of the domain, $\smash{\domainCentroid = 0}$, so that the particles verify $\smash{\sum_i \pvol_i\,\qq_i = 0}$ by \eqref{eq:partition-feasibility-centroid}.
Among the homotheties \eqref{eq:lifted-homothety} of the generators, we also fix their translation so that $\smash{\sum_i \pvol_i\,\dual_i = 0}$, which leaves the dilation as the only free scale of the normal.
%
On realizable particles, the centroids coincide with the particles, so the quantization cost reduces to
\begin{equation}
  \cpot_\infty(\qq) = \int_\domain \norm{\xx}^2 d\xx - \sum\nolimits_i \pvol_i \norm{\qq_i}^2 .
  \label{eq:flow-quantization}
\end{equation}
The motion of the particles is then governed by the Lagrangian
\begin{equation}
  \lag(\qq, \dot\qq) = \frac{1}{2} \sum\nolimits_i \pvol_i \norm{\dot\qq_i}^2 - \frac{1}{2\stiff^2}\, \cpot_\infty(\qq) ,
  \label{eq:flow-lagrangian}
\end{equation}
restricted to realizable particles, where $\stiff \!>\! 0$ sets the strength of the potential.
Its Euler--Lagrange equations with the holonomic constraint read
\begin{equation}
  \pvol_i\,\ddot\qq_i = \frac{1}{\stiff^2}\, \pvol_i\,\qq_i - s\,\pvol_i\,\dual_i ,
  \label{eq:flow-continuous}
\end{equation}
where the multiplier $s$ is the magnitude of the constraint force along the normal $\dual$, determined by the requirement that the particles remain realizable.
%
The potential thus pushes the particles outward, and the constraint absorbs the normal component of this push.
Its tangential component vanishes only when $\qq$ is parallel to $\dual$, \ie when the particles are themselves generators of their power diagram, up to a dilation.
The stationary configurations of the potential on the constraint are therefore the centroidal power diagrams with volumes $\pvol$, so the potential acts as a regularization of the cells, while $\stiff \!\to\! \infty$ leaves only the motion induced by the constraint.

\paragraph{Time discretization.}
We denote the time increment by $\tstep$, and use a superscript to indicate a time step, \eg $\smash{\qq^k}$ refers to the particle positions at time $\smash{t^k = k\tstep}$.
Following \citet{marsden2001discrete}, we discretize \eqref{eq:flow-lagrangian} by the discrete Lagrangian
\begin{equation}
  \lag_d(\qq^k, \qq^{k+1})
    = \frac{1}{2\tstep} \sum\nolimits_i \pvol_i \norm{\qq^{k+1}_i - \qq^k_i}^2 - \frac{\tstep}{4\stiff^2} \bigl( \cpot_\infty(\qq^k) + \cpot_\infty(\qq^{k+1}) \bigr) ,
  \label{eq:flow-discrete-lagrangian}
\end{equation}
which applies the trapezoidal rule to the potential.
Making the discrete action $\sum_k \lag_d(\qq^k, \qq^{k+1})$ stationary under the holonomic constraint gives the RATTLE integrator \cite{Andersen1983Rattle}, \ie the constrained version of the Störmer--Verlet scheme, which is symplectic and time-reversible \cite{marsden2001discrete, Hairer2006GNI}:
\begin{subequations}
  \label{eq:flow-rattle}
  \begin{empheq}[left={\empheqlbrace\,}]{align}
    \vel^{k+1/2}_i &= \vel^k_i + \frac{\tstep}{2} \Bigl( \frac{\qq^k_i}{\stiff^2} - s^k\,\dual^k_i \Bigr) , \label{eq:flow-rattle-kick} \\[1mm]
    \qq^{k+1}_i &= \qq^k_i + \tstep\, \vel^{k+1/2}_i , \label{eq:flow-rattle-drift} \\[1mm]
    \vel^{k+1}_i &= \vel^{k+1/2}_i + \frac{\tstep}{2} \Bigl( \frac{\qq^{k+1}_i}{\stiff^2} - r^{k+1}\,\dual^{k+1}_i \Bigr) . \label{eq:flow-rattle-kick2}
  \end{empheq}
\end{subequations}
Here, $\dual^k$ are the generators of the power diagram that maximizes $\jointdual$ with sites and target centroids at the particles $\qq^k$, the magnitude $s^k$ is chosen so that the new particles $\qq^{k+1}$ are realizable, and the magnitude $r^{k+1}$ is chosen so that the new velocities are tangent to the constraint, \ie $\smash{\sum_i \pvol_i\,(\dual^{k+1}_i)\tr \vel^{k+1}_i = 0}$.
The latter condition is linear in $r^{k+1}$ and is thus solved in closed form, reusing the generators $\dual^{k+1}$ required by the next step.
%
Eliminating the velocities yields the position recurrence
\begin{equation}
  \qq^{k+1}_i - 2\,\qq^k_i + \qq^{k-1}_i = \frac{\tstep^2}{\stiff^2}\,\qq^k_i - \tstep^2\, \bar{s}^k\, \dual^k_i ,
  \label{eq:flow-shake}
\end{equation}
with $\smash{\bar{s}^k = (s^k + r^k)/2}$, which is the SHAKE algorithm \cite{Ryckaert1977Shake}.
Other quadratures of the potential, such as its evaluation at $\qq^k$ only, which leads to the symplectic Euler scheme, give the same recurrence and only change how the velocities are staggered.
In all cases, the constraint force of a time step is directed along the generators $\dual^k$ of the \emph{current} power diagram.

\paragraph{Line search.}
By \eqref{eq:flow-rattle-kick} and \eqref{eq:flow-rattle-drift}, each time step moves the particles to
\begin{equation}
  \qq^{k+1} = \Bigl( 1 + \frac{\tstep^2}{2\stiff^2} \Bigr) \qq^k + \tstep\, \vel^k - \frac{\tstep^2}{2}\, s^k\, \dual^k ,
  \label{eq:flow-shake-step}
\end{equation}
\ie a prediction corrected along the generators $\dual^k$.
The direction $\dual^k$ requires solving the inverse problem of \S\ref{sec:joint-constraints} at the current particles, whose solution is unique only up to the homotheties of \eqref{eq:joint-nullspace}.
The translation is fixed by centering the generators, and the dilation is precisely the free magnitude $s^k$.
Unlike classical applications of SHAKE, where the constraint is an explicit function whose root yields this magnitude, realizable centroids are only described implicitly, as the exposed points of $\achievable$.
A naive approach thus requires a line search along $\dual^k$ until the new particles are realizable, each trial of which is a maximization of $\jointdual$.

\paragraph{Solver.}
Instead, the line search can be cast as a single concave maximization, which we solve by the proximal iteration of \S\ref{sec:regularized-joint}.
Writing the step \eqref{eq:flow-shake-step} as $\smash{\qq^{k+1} = \pred^{k+1} - s\,\dual^k - \translation}$, with the prediction $\smash{\pred^{k+1} = (1 + \tstep^2/(2\stiff^2))\,\qq^k + \tstep\,\vel^k}$, the unknown magnitude $\smash{s = \tstep^2 s^k / 2}$, and an unknown translation $\smash{\translation \in \R^d}$ that keeps the new particles centered, we maximize $\jointdual$ with sites $\qq^k$ and target centroids $\smash{\pred^{k+1}}$, subject to the linear constraints
\begin{equation}
  \sum\nolimits_i \pvol_i\,(\dual^k_i)\tr \shift_i = 0 , \qquad \sum\nolimits_i \pvol_i\,\shift_i = 0
  \label{eq:flow-line-constraint}
\end{equation}
on the shifts, which remove the shift components of the dilation and translation modes \eqref{eq:joint-nullspace}.
By the shift gradient \eqref{eq:joint-gradient-shift}, the multipliers of these constraints move the targets to $\smash{\pred^{k+1} - s\,\dual^k - \translation}$, and their values at the maximizer are the magnitude $s$ and the translation $\translation$.
Since the generators $\dual^k$ are centered, the identity \eqref{eq:partition-feasibility-centroid} at the maximizer fixes the translation in closed form, $\smash{\translation = \sum_i \pvol_i\,\pred^{k+1}_i / \domainVol}$, independently of $s$, which vanishes when the total momentum $\smash{\sum_i \pvol_i\,\vel^k_i}$ is zero.
More precisely, this problem is dual to the minimization of $s$ subject to $\smash{\pred^{k+1} - s\,\dual^k - \translation \!\in\! \achievable}$ over $(s, \translation)$ (Lemma~\ref{lem:line-constrained}), so it returns the smallest correction that brings the prediction back to $\achievable$, \ie the root of the line search closest to the prediction, without any bisection.
Moreover, the generators of its maximizer are a homothety of the generators of the new power diagram, so the direction $\smash{\dual^{k+1}}$ of the next step, and of the velocity step \eqref{eq:flow-rattle-kick2}, is a by-product of the same solve, and the inverse problem is only required to initialize the simulation, or not at all when it starts from a centroidal power diagram, with $\smash{\dual^0 = \qq^0}$.

As for $\jointdual$ in \S\ref{sec:regularized-joint}, this constrained maximization may have no solution, so we solve it by the proximal iteration \eqref{eq:regularized-proximal}, restricted to the constraints \eqref{eq:flow-line-constraint}.
Starting from the sites $\smash{\dual^{k,0} = \qq^k}$, each iteration maximizes $\jointdual_\alpha$ with sites $\smash{\dual^{k,m}}$, target centroids $\smash{\pred^{k+1}}$ and the constraints \eqref{eq:flow-line-constraint}, which only border its Newton system with $d \!+\! 1$ rows, and updates the generators as
\begin{equation}
  \dual^{k,m+1} = \dual^{k,m} + \alpha \bigl( \pred^{k+1} - s^{k,m+1}\,\dual^k - \translation^{k,m+1} - \centroid^{k,m+1} \bigr) ,
  \label{eq:flow-line-proximal}
\end{equation}
where $\smash{s^{k,m+1}}$ and $\smash{\translation^{k,m+1}}$ are the multipliers of the constraints and $\smash{\centroid^{k,m+1}}$ are the centroids of the resulting power diagram.
Each iteration is a strongly concave maximization, which always has a solution, and the iterates converge for any $\alpha \!>\! 0$ to the solution of the line search, with the new particles as centroids and a homothety of their generators as sites (Lemma~\ref{lem:line-constrained}).
%
The remaining limitations concern the time step rather than the solver.
If the line misses $\achievable$, which may happen for large steps, the line search has no solution: the generators then grow without bound, which is detected, while the centroids converge to the achievable centroids closest to the line, as in Lemma~\ref{lem:proximal-any}.
If the line enters $\achievable$ through a flat face, the new particles are achievable but not realizable, and the generators collapse.
Both failures are excluded for small enough time steps, since the prediction then lies close to $\qq^k$, where $\dual^k$ is normal to a smooth part of the boundary of $\achievable$.
The holonomic formulation is thus variational and solvable, but each time step costs several maximizations of $\jointdual_\alpha$, and the time step is limited by these failures, which motivates the relaxations below.
% TODO: numerical validation of the solver (iterations per step, step limit) in ../src.

\subsection{Relaxation by Penalty}
\label{sec:flow-softened}

The simplest relaxation removes the holonomic constraint by using the particles as the \emph{sites} of the power diagram rather than as its centroids, and replaces the centroid constraint by the penalty of \S\ref{sec:regularized-joint}.
Using $\jointdual_\alpha$ in place of $\jointdual$ in \eqref{eq:flow-centroid-potential}, again with sites and target centroids at the particles, we define the potential
\begin{equation}
  \cpot_\alpha(\qq) = \max_{\shift, \weight} \jointdual_\alpha(\shift, \weight) ,
  \label{eq:flow-softened-potential}
\end{equation}
which is finite for any distinct particles, so the Lagrangian \eqref{eq:flow-lagrangian} with $\cpot_\alpha$ in place of $\cpot_\infty$ is unconstrained.
By \eqref{eq:regularized-projection}, the centroids $\centroid$ of the maximizer are the projection onto $\achievable$ of the dilated particles $(1 + 1/\alpha)\,\qq$, which is the dilate-and-project construction of \citet{Bourne2026ConvexOrder}, here used as a potential.
By \eqref{eq:regularized-gradient-shift}, the optimal shifts are $\smash{\shift_i = \alpha\,(\qq_i - \centroid_i)}$, so the particles are no longer the centroids of their cells, but lie at a distance $\smash{\norm{\shift_i}/\alpha}$ from them.

To derive the forces, note that the particles enter $\jointdual_\alpha$ both as sites and as targets, while the penalty in \eqref{eq:regularized-dual} does not depend on them.
At fixed multipliers, the sites contribute the residual \eqref{eq:centroid-gradient} and the targets contribute the derivative of the centroid constraint term in \eqref{eq:joint-lagrangian}, so that, by the envelope theorem and since the volumes are exact,
\begin{equation}
  \partial_{\qq_i} \cpot_\alpha
  = 2\,\pvol_i\,(\qq_i - \centroid_i) + 2\,\pvol_i\,\shift_i
  = 2\,(1 + \alpha)\,\pvol_i\,(\qq_i - \centroid_i) .
  \label{eq:softened-site-gradient}
\end{equation}
The Störmer--Verlet scheme for this unconstrained Lagrangian then reads
\begin{subequations}
  \label{eq:flow-softened-step}
  \begin{empheq}[left={\empheqlbrace\,}]{align}
    \vel^{k+1/2}_i &= \vel^k_i + (1 + \alpha)\, \frac{\tstep}{2\stiff^2} \bigl( \centroid^k_i - \qq^k_i \bigr) , \label{eq:flow-softened-step-kick} \\[1mm]
    \qq^{k+1}_i &= \qq^k_i + \tstep\, \vel^{k+1/2}_i , \label{eq:flow-softened-step-drift} \\[1mm]
    \vel^{k+1}_i &= \vel^{k+1/2}_i + (1 + \alpha)\, \frac{\tstep}{2\stiff^2} \bigl( \centroid^{k+1}_i - \qq^{k+1}_i \bigr) , \label{eq:flow-softened-step-kick2}
  \end{empheq}
\end{subequations}
where $\centroid^k$ are the centroids of the power diagram that maximizes $\jointdual_\alpha$ with sites and target centroids at $\qq^k$.
Since the centroids $\centroid^{k+1}$ are reused in the next step, each step costs a single maximization of $\jointdual_\alpha$.
%
This update is explicit, symplectic and time-reversible, and, contrary to the holonomic formulation, always has a solution.
The price is that the constraint is only approximated: the cells keep their prescribed volumes exactly, but their centroids differ from the particles by $\smash{\shift_i/\alpha}$.
In the limit $\alpha \!\to\! \infty$, $\jointdual_\alpha$ recovers $\jointdual$ (\S\ref{sec:regularized-joint}), and the potential $\cpot_\alpha$ recovers the holonomic formulation.
% TODO: step-size restriction of this explicit scheme as a function of alpha (force scaled by 1 + alpha, acting on a gap of order 1/alpha).

\subsection{Relaxation by Projection}
\label{sec:flow-projection}

Alternatively, we can keep the holonomic constraint exact and relax the integrator instead.
To this end, we evaluate both the potential force and the constraint force of \eqref{eq:flow-continuous} at the \emph{end} of the time step, \ie by an implicit Euler step:
\begin{equation}
  \qq^{k+1}_i - 2\,\qq^k_i + \qq^{k-1}_i = \frac{\tstep^2}{\stiff^2}\,\qq^{k+1}_i - \tstep^2\, s\, \dual^{k+1}_i ,
  \label{eq:flow-implicit}
\end{equation}
which differs from \eqref{eq:flow-shake} in that the generators are those of the \emph{new} power diagram.
With the velocities $\smash{\vel^k = (\qq^k - \qq^{k-1})/\tstep}$ and for $\tstep \!<\! \stiff$, rearranging \eqref{eq:flow-implicit} gives
\begin{equation}
  \qq^{k+1} = \pred^{k+1} - s'\,\dual^{k+1} ,
  \qquad
  \pred^{k+1} = (1 + \predcoef) \bigl( \qq^k + \tstep\, \vel^k \bigr) ,
  \label{eq:flow-projection-pred}
\end{equation}
with the coefficient $\smash{\predcoef = \tstep^2 / (\stiff^2 - \tstep^2)}$ and the magnitude $\smash{s' = (1 + \predcoef)\,\tstep^2 s}$.
The prediction $\smash{\pred^{k+1}}$ is thus the free drift of the particles, dilated by the implicit treatment of the potential.
%
Since the generators $\smash{\dual^{k+1}}$ are an outward normal of $\achievable$ at the centroids $\smash{\qq^{k+1}}$ of their power diagram (\S\ref{sec:background}), the correction $\smash{\pred^{k+1} - \qq^{k+1} = s'\,\dual^{k+1}}$ with $s' \!>\! 0$ lies in the normal cone of $\achievable$ at $\smash{\qq^{k+1}}$.
By the classical characterization of projections onto convex sets \cite[Thm.~3.16]{Bauschke2017Convex}, this is exactly the condition for $\smash{\qq^{k+1}}$ to be the point of $\achievable$ closest to $\smash{\pred^{k+1}}$ in the volume-weighted $\smash{L^2}$ sense.
Equivalently, $\smash{\qq^{k+1}}$ minimizes over $\achievable$ the incremental potential
\begin{equation}
  \frac{1}{2} \sum\nolimits_i \pvol_i \norm{\qq_i - \qq^k_i - \tstep\,\vel^k_i}^2 - \frac{\tstep^2}{2\stiff^2} \sum\nolimits_i \pvol_i \norm{\qq_i}^2 ,
  \label{eq:flow-incremental}
\end{equation}
whose second term is the potential \eqref{eq:flow-quantization} up to a constant, and which is strongly convex for $\tstep \!<\! \stiff$.
The line search of \S\ref{sec:flow-holonomic} along the old normal $\dual^k$ is thus replaced by a projection along the new normal $\smash{\dual^{k+1}}$, which always exists.

The coefficient $\predcoef$ unifies two integrators.
For $\smash{\predcoef = \tstep^2/(\stiff^2 - \tstep^2)}$, the update is the implicit Euler step \eqref{eq:flow-implicit} of the constrained dynamics.
For $\predcoef \!=\! 0$, \ie without potential ($\stiff \!\to\! \infty$), the update projects the free drift onto the constraint, as in the projection methods of geometric integration \cite{Hairer2006GNI} and their counterparts in graphics \cite{Goldenthal2007FastProjection, Muller2007PBD}.
Projections onto a convex set at every step also underlie the catching-up algorithm of \citet{Moreau1977Sweeping} and its time-stepping extensions to unilateral contact \cite{Jean1999NSCD}.
%
Two results make this relaxation work in our setting: the prediction is never achievable, and the projection is computed by the softened dual $\jointdual_\alpha$.

\paragraph{Unachievable predictions.}
If the prediction were achievable, it would be its own projection, with a vanishing normal and no power diagram.
This never happens.
Indeed, since the previous particles $\qq^{k-1}$ are achievable, the velocities have a non-negative component along the generators $\dual^k$ of the current power diagram, \ie $\smash{\sum_i \pvol_i\,(\dual^k_i)\tr \vel^k_i \ge 0}$ (Lemma~\ref{lem:normal-velocity}).
Under this condition, the prediction \eqref{eq:flow-projection-pred} is unachievable for any $\predcoef \!>\! 0$, and for $\predcoef \!=\! 0$ unless the particles are at rest, in which case they are correctly a fixed point (Corollary~\ref{cor:unachievable-prediction}).
For $\predcoef \!>\! 0$, the dilation pushes the prediction strictly beyond the supporting hyperplane of $\achievable$ at $\qq^k$ with normal $\dual^k$, by an amount independent of the velocities.
%
These results hold at every step $k \!\ge\! 1$, so that only the initial velocities require attention.
In the common case where the simulation starts from a centroidal power diagram, \ie $\smash{\dual^0 = \qq^0}$, Lemma~\ref{lem:unachievable-prediction} requires
\begin{equation}
  (1 + \predcoef)\, \tstep \sum\nolimits_i \pvol_i\,(\qq^0_i)\tr \vel^0_i + \predcoef \sum\nolimits_i \pvol_i \norm{\qq^0_i}^2 \ge 0 ,
  \label{eq:flow-initial-velocity}
\end{equation}
whose first term is proportional to the rate of change of the spread $\smash{\tfrac12 \sum_i \pvol_i \norm{\qq_i}^2}$ of the particles about $\domainCentroid$.
For initial velocities given by the cell averages of a divergence-free field tangent to $\partial\domain$, this first term vanishes, since the flow keeps the cells in $\achievable$ for small times of either sign and $\qq^0$ maximizes the spread there.
Otherwise, the inward normal component can be removed by replacing $\vel^0$ with $\smash{\vel^0 - \min(0, \gamma)\,\qq^0}$, where $\smash{\gamma = \sum_i \pvol_i\,(\qq^0_i)\tr \vel^0_i / \sum_i \pvol_i \norm{\qq^0_i}^2}$, which preserves the zero total momentum.

\paragraph{Projection by the softened dual.}
By \eqref{eq:regularized-projection}, maximizing $\jointdual_\alpha$ with sites $\qq$ and target centroids $\tcentroid$ returns the power diagram whose centroids are the projection onto $\achievable$ of the pushed targets $\smash{\tcentroid + \qq/\alpha}$.
We thus set the sites to the drift $\smash{\qq^k + \tstep\,\vel^k}$, and the target centroids to $\smash{\tcentroid^{k+1} = \pred^{k+1} - (\qq^k + \tstep\,\vel^k)/\alpha}$, so that the pushed targets coincide with the prediction.
The new particles then do not depend on $\alpha$, which only scales the generators $\smash{\dual^{k+1} = \alpha\,(\pred^{k+1} - \qq^{k+1})}$ given by \eqref{eq:regularized-gradient-shift}.
In fact, the softened problems with stiffness $\alpha$ and $\lambda\alpha$ coincide up to the scaling $\dual \mapsto \lambda\,\dual$ and a factor $\lambda$ on the objective (Lemma~\ref{lem:stiffness-scaling}), so they share the same conditioning.
For $\predcoef \!>\! 0$, the choice $\alpha = 1/\predcoef$ places the targets at the drift, and, for $\predcoef \!=\! 0$, any $\alpha \!>\! 0$ applies, with targets $\smash{(1 - 1/\alpha)(\qq^k + \tstep\,\vel^k)}$.
In both cases, the initial guess $\shift = 0$ places the generators at the drift.
The relaxed update then takes the form
\begin{subequations}
  \label{eq:flow-projection-step}
  \begin{empheq}[left={\empheqlbrace\,}]{align}
    \pred^{k+1} &= (1 + \predcoef) \bigl( \qq^k + \tstep\, \vel^k \bigr) , \label{eq:flow-projection-step-pred} \\[1mm]
    \qq^{k+1} &= \pred^{k+1} - \dual^{k+1} / \alpha , \label{eq:flow-projection-step-pos} \\[1mm]
    \vel^{k+1} &= \bigl( \qq^{k+1} - \qq^k \bigr) / \tstep , \label{eq:flow-projection-step-vel}
  \end{empheq}
\end{subequations}
where $\smash{\dual^{k+1}}$ are the generators of the power diagram that maximizes $\jointdual_\alpha$ with the sites and targets above, and whose centroids are the new particles $\smash{\qq^{k+1}}$.

\paragraph{Explicit treatment of the potential.}
Since the sites do not affect the projection, they only determine where the potential is evaluated.
Setting the sites to the current particles $\qq^k$ instead, with targets at the drift and $\smash{\alpha = \stiff^2/\tstep^2}$, the pushed targets become $\smash{(1 + \tstep^2/\stiff^2)\,\qq^k + \tstep\,\vel^k}$, which treats the potential explicitly and the constraint implicitly, as a forward--backward step.
This variant has no restriction on the time step, and its prediction is also unachievable by Corollary~\ref{cor:unachievable-prediction}.
The explicit or implicit treatment of the potential thus amounts to placing the sites of the same maximization at the current particles or at the drift.

\paragraph{Properties.}
Since $\achievable$ is closed and convex, the projection always exists and is unique for any prediction and any time step, and it is computed by a single concave maximization rather than a line search.
All particles remain in $\achievable$ and thus satisfy the conservation identity \eqref{eq:partition-feasibility-centroid} at every step, so the total momentum $\smash{\sum_i \pvol_i\,\vel^{k+1}_i}$ vanishes exactly.
For $\predcoef \!=\! 0$, the nonexpansiveness of projections onto convex sets, together with $\smash{\qq^k = \proj(\qq^k)}$, gives $\smash{\sum_i \pvol_i \norm{\vel^{k+1}_i}^2 \le \sum_i \pvol_i \norm{\vel^k_i}^2}$, so the kinetic energy never increases, and the energy lost is the normal velocity removed by the projection, as in an inelastic impact.
These properties come at the price of the variational structure of \S\ref{sec:flow-holonomic}: since the projection is not invertible, the relaxed update is neither symplectic nor time-reversible, for any quadrature of the potential.
%
Finally, the projection is realizable only if the generators $\smash{\dual^{k+1}_i = \alpha\,(\pred^{k+1}_i - \qq^{k+1}_i)}$ are pairwise distinct, since the relative boundary of $\achievable$ also contains flat faces exposed by coinciding generators (\S\ref{sec:background}).
This is the case for small enough time steps, since the normals at boundary points near the realizable particles $\qq^k$ remain close to $\dual^k$, while the nonexpansiveness of the projection keeps $\smash{\qq^{k+1}}$ within a distance of order $\tstep$ of $\qq^k$.
When it fails, the projection still returns achievable centroids instead of having no solution.

\subsection{Relation to Prior Work}
\label{sec:flow-prior}

\paragraph{\citet{Gallouet2018Lagrangian}.}
For $\alpha \!=\! 0$, the shifts vanish and the penalty potential \eqref{eq:flow-softened-potential} reduces to the cost of a volume-constrained power diagram (\S\ref{sec:background}),
\begin{equation}
  \cpot_0(\qq) = \max_\weight \dualenergy(\weight) ,
  \label{eq:flow-gm-potential}
\end{equation}
which, with the factor $1/(2\stiff^2)$ of \eqref{eq:flow-lagrangian}, is the potential of \citet{Gallouet2018Lagrangian}.
The force \eqref{eq:softened-site-gradient} then pulls each particle toward the centroid of its own volume-constrained cell, and the symplectic Euler scheme of this Lagrangian, $\smash{\vel^{k+1}_i = \vel^k_i + \tfrac{\tstep}{\stiff^2} (\centroid^k_i - \qq^k_i)}$ and $\smash{\qq^{k+1}_i = \qq^k_i + \tstep\,\vel^{k+1}_i}$, is exactly their scheme \cite[Eq.~(1.11)]{Gallouet2018Lagrangian}.
They show its convergence towards smooth solutions of the Euler equations for $\tstep \!\le\! \stiff$, with an error bound that vanishes only for $\tstep \!=\! o(\stiff^2)$ \cite[Thm.~1.4]{Gallouet2018Lagrangian}, and they also report a better conservation of the energy with the velocity Verlet scheme \cite[Fig.~4]{Gallouet2018Lagrangian}, which is our update \eqref{eq:flow-softened-step} for $\alpha \!=\! 0$.
Their scheme is thus the $\alpha \!=\! 0$ member of the penalty relaxation of \S\ref{sec:flow-softened}: the cells always match their prescribed volumes, but the particles are tied to their centroids only through the strength $\stiff$ of the potential, so the velocities are divergence-free only in the limit $\stiff \!\to\! 0$, whereas $\alpha \!>\! 0$ adds a penalty that ties the particles to the centroids directly.
%
Their construction also explains the holonomic formulation of \S\ref{sec:flow-holonomic}.
Incompressible Euler flows are geodesics on the group of measure-preserving maps, and \citet{Gallouet2018Lagrangian} discretize such maps by their averages on a fixed partition of the domain into cells of volumes $\pvol$, penalizing the distance of these averages to the measure-preserving maps.
These averages range exactly over the achievable centroids $\achievable$ \cite{Bourne2026ConvexOrder}, so the holonomic formulation, which restricts them to the realizable ones, replaces this penalty by an exact constraint.

\paragraph{Power Particles.}
Power Particles \cite{deGoes2015PowerParticles} store, for each particle, a site $\qq_i$ and a velocity $\vel_i$, and advance them by an operator splitting.
Given sites $\qq^k$ and velocities $\vel^k$, and up to a cyclic rotation of \cite[Alg.~1]{deGoes2015PowerParticles}, each time step performs:
\begin{enumerate}
  \item \emph{Volume enforcement:} compute the weights that maximize $\dualenergy(\weight)$ for the sites $\qq^k$, and the centroids $\centroid^k$ of the resulting cells.
  \item \emph{Pressure projection:} solve $\smash{\tstep\,\Lmat\,\press = \Dmat\,\vel^k}$ and set $\smash{\vel^{k+1} = \vel^k - \tstep\,\Mmat^{-1} \Gmat\,\press}$, where $\smash{\Mmat = \mathrm{diag}(\pvol_i \Id_d)}$ is the diagonal mass matrix, $\Dmat$ is the discrete divergence of the volumes with respect to the sites at fixed weights, $\smash{\Gmat = -\Dmat\tr}$, and $\smash{\Lmat = \Dmat\,\Mmat^{-1} \Gmat}$ \cite[Eqs.~(4) and~(7)]{deGoes2015PowerParticles}.
  \item \emph{Advection:} $\smash{\qq^{k+1} = \centroid^k + \tstep\,\vel^{k+1}}$.
\end{enumerate}
This splitting is not derived from a variational principle, and \citet[Remark~1.5]{Gallouet2018Lagrangian} note that it lacks a Hamiltonian structure, although it behaves better than their penalty scheme in practice and does not depend on a stiffness.
Steps~1 and~2 resemble the two stages of RATTLE, but for the volume constraints: the weight solve enforces the volumes exactly on the advected sites, as a position step, and the pressure projection removes the velocity component that changes the volumes to first order, as a velocity step.

In terms of our formulations, the centroids $\centroid^k$ of step~1 are those of the volume-constrained power diagram generated by the advected sites, \ie the exposed point of $\achievable$ with normal $\smash{\centroid^{k-1} + \tstep\,\vel^k}$.
Since the projection of $\smash{(1 + \predcoef)\,\yy}$ onto $\achievable$ tends to the exposed point with normal $\yy$ as $\predcoef \!\to\! \infty$, the position update of Power Particles is the limit $\predcoef \!\to\! \infty$ of the projection relaxation of \S\ref{sec:flow-projection}, in which the drift is used as a normal rather than as the point to project.
%
The pressure projection, in turn, is a modified Newton step on $\jointdual$.
Denoting the negated Hessian blocks of $\jointdual$ by $\smash{\Kmat = -\partial^2_{\shift\shift}\jointdual}$ and $\smash{\wlap = -\partial^2_{\weight\weight}\jointdual}$, and the mixed block \eqref{eq:joint-hessian} by $\smash{\Bmat = \partial^2_{\shift\weight}\jointdual = -\partial\cellvol/\partial\shift}$, a Newton step from zero shifts solves
\begin{equation}
  \begin{bmatrix} \Kmat & -\Bmat \\ -\Bmat\tr & \wlap \end{bmatrix}
  \begin{bmatrix} \delta\shift \\ \delta\weight \end{bmatrix}
  =
  \begin{bmatrix} \partial_{\shift}\jointdual \\ \partial_{\weight}\jointdual \end{bmatrix} .
  \label{eq:flow-newton}
\end{equation}
Power Particles modify this step in three ways: they solve for the weights to convergence first, so that $\smash{\partial_{\weight}\jointdual = 0}$; they zero out the weight--weight block $\wlap$; and they replace the shift--shift block by the diagonal $\smash{\Kmat \approx 2\,\Mmat}$.
Note that, by \eqref{eq:centroid-hessian}, the block $\Kmat$ of $\jointdual$ consists of facet terms only, since the mass terms of the energy Hessian cancel with those of the shifted weights (Appendix~\ref{sec:appendix-centroid}), so this last modification is a lumping by a mass matrix, which coincides with the penalty block $\smash{\tfrac{2}{\alpha}\Mmat}$ of $\jointdual_\alpha$ for $\alpha \!=\! 1$ once the facet terms are dropped.
% TODO: confirm with the user how to phrase the "lumping" of the shift--shift block (row-sum lumping of the facet block does not give the mass matrix).
The modified system is the optimality system of a projection: its second row requires $\smash{\Bmat\tr \delta\shift = 0}$, \ie shifts that preserve the volumes to first order, and eliminating the shifts gives
\begin{equation}
  \bigl( \Bmat\tr \Mmat^{-1} \Bmat \bigr)\, \delta\weight = -\Bmat\tr \Mmat^{-1} \partial_{\shift}\jointdual ,
  \qquad
  \delta\shift = \tfrac12\,\Mmat^{-1} \bigl( \partial_{\shift}\jointdual + \Bmat\,\delta\weight \bigr) ,
  \label{eq:flow-pp-schur}
\end{equation}
where zeroing $\wlap$ is precisely what reduces the Schur complement to the pressure Laplacian $\smash{\Bmat\tr \Mmat^{-1} \Bmat}$.
With target centroids $\smash{\tcentroid = \centroid^k + \tstep\,\vel^k}$, the shift gradient \eqref{eq:joint-gradient-shift} at the converged weights is $\smash{\partial_{\shift}\jointdual = 2\,\Mmat\,\tstep\,\vel^k}$, so that $\smash{\delta\shift = \tstep\,\vel^{k+1}}$, where $\smash{\vel^{k+1}}$ is the mass-weighted projection of $\smash{\vel^k}$ onto the velocities that preserve the volumes to first order, with the weight update $\delta\weight$ playing the role of the pressure.
This is the pressure projection of step~2, up to the divergence: by \eqref{eq:shifted-site}, the shifted weights vary with the shifts as $\smash{2\,(\qq_i - \tcentroid_i)}$ at zero shifts, so $\smash{-\Bmat\tr}$ coincides with the divergence $\Dmat$ with respect to the sites only when the targets coincide with the sites.
% TODO: check constants and signs between delta w and dt*pi against [deGoes2015PowerParticles, Eqs. (4) and (7)] (pp.pdf not readable in the sandbox).
%
Power Particles can thus be read as a RATTLE-like splitting for the volume constraints, whose position step is the $\predcoef \!\to\! \infty$ limit of our projection and whose velocity step is a single modified Newton step on $\jointdual$.
In contrast with the velocity step \eqref{eq:flow-rattle-kick2} of the holonomic formulation, which removes the single normal component along $\smash{\dual^{k+1}}$, the pressure projection removes one component per cell, since it linearizes the $n$ volume constraints rather than the single realizability constraint.

\end{document}
